
\Problem{Guanyadors}

\Statement

Suposant la declaració de la tupla \verb|Player| següent
\begin{verbatim}
struct Player { 
    string name; 
    int score;
};
\end{verbatim}
fes una funció amb la capçalera
\begin{verbatim}
vector<string> winners(vector<Player>::iterator begin,
                       vector<Player>::iterator end);
\end{verbatim}
que retorni un vector amb els jugadors amb la puntuació
màxima. Tingues en compte que pot haver-hi empat múltiple.

\Input

La funció rep dos iteradors indicant principi i final d'una subseqüència en un vector de jugadors. Aquesta subseqüència pot tenir qualsevol longitud, fins i tot ser buida. La subseqüència també podria ser el vector sencer.

\Output

La funció retorna un vector amb els noms dels jugadors que tenen
alhora la puntuació màxima en el mateix ordre que es troben al vector
original.

\Observation

Als fitxers públics (icona del gatet) trobaràs: \verb|main.cc|, \verb|player.hh| (declaració de la tupla \verb|Player|), \verb|guanyadors.cc| (a on has d'implementar la funció). També hi ha un \verb|Makefile| i el directori \verb|.vscode| que té la configuració per compilar i debuggar amb VSCode.

Has d'implementar \verb|winners| en el fitxer \verb|guanyadors.cc|, compilar (està preparat per poder compilar i debuggar amb VSCode), i enviar només el fitxer \verb|guanyadors.cc|.

\SampleTwoCol
